\section{A Geometric Direction of Repair Success}
\label{sec:repair_geometry}

When the model is given the chance to repair its own failing code, the question is whether the shift in its hidden state between the failed attempt and the repair attempt carries a stable, decodable signature of whether the repair will succeed. The headline answer of this section is that a statistically detectable geometric direction does exist in the raw deltas, and is statistically significant on two complementary tests, but does not survive a conditional residualization once observable repair-context covariates are checked and found to differ between the success and failure groups.

For each task where attempt 0 fails the unit tests and attempt 1 exists, the per-task hidden state delta is defined as
\[
    \delta_i = h_{\text{layer 29}}(\text{attempt 1}_i) - h_{\text{layer 29}}(\text{attempt 0}_i),
\]
and the candidate contrastive direction is
\[
    v = \mathrm{mean}\bigl(\{\delta_i : i \in \text{success}\}\bigr) - \mathrm{mean}\bigl(\{\delta_i : i \in \text{failure}\}\bigr),
\]
where success and failure refer to whether attempt 1 passes or fails the unit tests. Each attempt is a fresh forward pass on the full prompt with no state carried over from previous attempts, as described in Section~\ref{sec:setup}, so $\delta_i$ is driven entirely by what the model re-reads on attempt 1: the failing code from attempt 0, the captured error output, and the retry instruction.

\paragraph{Data cleaning.}
Two filters are applied before computing deltas, both targeting the same underlying pathology, model generations that hit the 768-token cap with repetitive broken code, at two different scopes: trajectory-wide and attempt-0 only. The first filter removes the 47 tasks tagged as repetition-loop tasks, where the model is stuck producing broken output that hits the 768-token cap across all five attempts so the entire trajectory is contaminated. After this initial pass, 327 candidate 0$\to$1 transitions remain. Of these, a further 91 are excluded because attempt 0 alone hit the same cap: even when attempt 1 produced valid code, the truncated attempt-0 generation is pasted verbatim into the repair prompt and contaminates the delta for the failure group, with 83 of the 91 failing at attempt 1. The final clean sample contains 236 transitions, 128 with a successful repair and 108 with a failed repair.

\subsection{Raw tests}
\label{subsec:raw_tests}

Two complementary tests are run on the contrastive direction $v$ in the full 2560-dimensional layer-29 space. Both compare the real statistic to an empirical null distribution generated by repeatedly shuffling the pass/fail labels across the 236 deltas.

The first test, on magnitude, asks whether the success/failure mean difference is larger than label noise can explain. The norm $\|v\|$ is computed on the real assignment and on 1000 random shufflings of the labels. The second test, on inclination, asks whether the direction is geometrically stable: the 236 deltas are split into two stratified halves (each preserving the 128/108 success/failure ratio), $v$ is recomputed independently on each half, and the cosine similarity between the two recovered directions is taken. This is repeated 1000 times. A null distribution is generated by the same split-half procedure on shuffled-label data. Results for both tests are reported in Table~\ref{tab:raw_direction_tests}.

\begin{table}[h]
\centering
\begin{tabular}{lccc}
\toprule
Test                  & Real value & Null mean & p-value  \\
\midrule
Magnitude $\|v\|$     & $9.76$     & $3.64$    & $< 0.001$ \\
Split-half cosine     & $0.767$    & $-0.011$  & $< 0.001$ \\
\bottomrule
\end{tabular}
\caption{Two complementary tests of the contrastive direction $v$ on the 236 clean $0 \to 1$ transitions. The magnitude null has std $0.59$ and max $6.58$; the real magnitude is therefore roughly ten standard deviations above the null mean, and 0 of 1000 shuffles reached it. The split-half cosine null has mean $-0.011$ and std $0.225$; the real mean cosine of $0.767$ is therefore over three standard deviations above the null mean.}
\label{tab:raw_direction_tests}
\end{table}

Taken on their own, these two tests look like clean evidence for a repair-relevant feature: the magnitude sits roughly ten standard deviations above the null mean and is larger than every one of the 1000 shuffled norms, and the inclination is moderately stable across random data splits. The qualifier is introduced in the next subsection.

\FloatBarrier
\subsection{Sensitivity check}
\label{subsec:sensitivity_check}

The contrastive construction $\mathrm{mean}(\text{success}) - \mathrm{mean}(\text{failure})$ is a difference of group means, so it already cancels the linear contribution of any covariate whose mean is equal across the two groups. For a one-hot-encoded categorical covariate this condition is equivalent to its category proportions being equal across groups, that is, the same categorical distribution. Residualization is therefore necessary only when a candidate covariate's group means differ. Otherwise the operation removes within-group variance without changing between-group means and collapses the signal artifactually. The principled procedure is to test, before residualizing, whether each candidate covariate differs between the success and failure groups, and to include in the residualization only those that do.

Three candidate covariates are computed per delta from the attempt-level log:
\begin{itemize}
    \item $\Delta\texttt{prompt\_tokens}$, the change in input token length between attempt 1 and attempt 0.
    \item $\texttt{code\_length\_attempt\_0}$, the number of completion tokens produced at attempt 0, used as a proxy for how much wrong code gets pasted into the repair prompt.
    \item $\texttt{failure\_type}$ of attempt 0, a categorical with three observed levels (assertion, runtime, syntax). Of the five non-pass labels in the taxonomy described in Section~\ref{sec:setup}, only assertion, runtime, and syntax appear in the clean sample; timeout and the unknown fallback are absent.
\end{itemize}

Welch's two-sample t-test, which does not assume equal variance, is used for the two continuous covariates and a chi-squared test on the $3 \times 2$ contingency table is used for the categorical covariate. Each test produces a p-value reported in Table~\ref{tab:covariate_check}. The null hypothesis is that the success and failure group means are equal for the two continuous covariates, and that the failure-type proportions are equal across the two groups for the categorical covariate.

\begin{table}[h]
\centering
\begin{tabular}{lccc}
\toprule
Covariate & Success mean (n=128) & Failure mean (n=108) & p-value \\
\midrule
$\Delta\texttt{prompt\_tokens}$              & $337.1$ & $431.5$ & $< 0.001$ \\
$\texttt{code\_length\_attempt\_0}$          & $188.8$ & $290.9$ & $< 0.001$ \\
$\texttt{failure\_type}$ of attempt 0        & \multicolumn{2}{c}{(categorical, $3$ levels)} & $< 0.001$ \\
\bottomrule
\end{tabular}
\caption{Covariate balance check on the 236 clean 0$\to$1 transitions. All three observable covariates differ significantly between the success and failure groups, which warrants residualizing against them.}
\label{tab:covariate_check}
\end{table}

All three covariates differ at $p < 0.001$. Successfully repaired tasks tend to come from shorter attempt-0 code, from a shorter prompt growth, and from a different mix of failure types (more runtime errors, fewer assertion errors). The contrastive subtraction cancels only the linear contribution of covariates whose group means are equal. Since the group means of all three differ, the subtraction does not cancel them, and residualization against them is conceptually warranted.

The operation is now described formally. Let $X \in \mathbb{R}^{236 \times 4}$ denote the covariate matrix with one row per task and four columns (the two continuous covariates and the two failure-type dummies, with assertion as the baseline level). For each of the 2560 hidden state dimensions $d$, an OLS regression is fit on the 236 deltas with the $d$-th coordinate as response and $X$ as predictor matrix, producing an intercept $\hat{\alpha}^{(d)} \in \mathbb{R}$ and a slope vector $\hat{\beta}^{(d)} \in \mathbb{R}^{4}$. The residualized delta for task $i$ in dimension $d$ is then $\tilde\delta_{i,d} = \delta_{i,d} - \bigl(\hat{\alpha}^{(d)} + X_i \hat{\beta}^{(d)}\bigr)$, where $X_i$ is the covariate row for task $i$. Stacking across all 2560 dimensions yields a residualized delta matrix of the same shape as the original. The contrastive direction $v$ is recomputed on the residualized deltas, and the two tests are rerun. Table~\ref{tab:residualization} reports the raw and residualized statistics side by side.

\begin{table}[h]
\centering
\begin{tabular}{lcc}
\toprule
Statistic & Raw & Residualized \\
\midrule
Magnitude $\|v\|$             & $9.76$   & $3.72$ \\
Magnitude p-value             & $< 0.001$ & $0.126$ \\
Split-half cosine             & $0.767$  & $0.140$ \\
Split-half p-value            & $< 0.001$ & $0.942$ \\
\bottomrule
\end{tabular}
\caption{Raw versus conditional residualization on the 236 clean $0 \to 1$ transitions. After residualizing against the three observable covariates that differ between groups, neither the magnitude test nor the split half test reaches significance.}
\label{tab:residualization}
\end{table}

The signal does not survive the full conditional residualization. The residualized norm of $3.72$ retains only $38.1\%$ of the raw norm $9.76$, the split-half cosine drops from $0.767$ to $0.140$, and neither the magnitude test ($p = 0.126$) nor the split-half test ($p = 0.942$) reaches significance. This should not be read as the raw signal being false: the raw direction passes both statistical tests by wide margins. The honest reading is that the observed geometric direction is a correlate of repair success, substantially mediated by what differs in the repair context, rather than evidence of an isolated repair-comprehension direction in activation space.

Two of the three covariates used in the residualization are partially redundant: $\Delta\texttt{prompt\_tokens}$ decomposes into the failed code length, the error output length, and an instruction template that is constant up to a single \texttt{failure\_type} token (the \texttt{failure\_type} variation is one of the three covariates already included in the residualization matrix, so its contribution is subtracted explicitly), so the failed code length is also present in $\texttt{code\_length\_attempt\_0}$. Empirically, the two are nearly collinear, with a Pearson $r$ of $0.995$. An extra regression column, however, absorbs on average about $1/n$ of the variance even when it carries no real information, and the redundant $\Delta\texttt{prompt\_tokens}$ column is a single parameter against 236 deltas, so its spurious contribution to the in-sample $R^2$ of the per-dimension OLS fit is about 0.4 percentage points. The measured value agrees: adding the redundant column on top of the other three raises the mean in-sample $R^2$ by under one percentage point. The residualization therefore removes essentially the same variance whether the redundant column is included or not, and any overfitting bias on the residualized norm is negligible.

\subsection{Beyond the first attempt}

The magnitude and inclination tests above are estimated on the $0 \to 1$ transition only, where 128 successes and 108 failures provide enough data for the contrastive construction. At later transitions the situation changes sharply, as Table~\ref{tab:pass_rates} shows.

\begin{table}[h]
\centering
\begin{tabular}{cc}
\toprule
Attempt & Pass rate \\
\midrule
0 & 15.8\% \\
1 & 36.4\% \\
2 & 6.7\%  \\
3 & 2.7\%  \\
4 & 1.4\%  \\
\bottomrule
\end{tabular}
\caption{Per-attempt pass rates among tasks that reached each attempt (the denominator shrinks across rows as tasks exit the trajectory after passing). The bulk of the gain happens at attempt 1; later attempts contribute only marginally.}
\label{tab:pass_rates}
\end{table}

The first piece of error feedback delivers most of the information the model can use, and subsequent retries yield diminishing returns. This is consistent with prior reports that iterative self-repair plateaus after the first one or two attempts \citep{arimbur2026howmanytries}. At transitions $1 \to 2$, $2 \to 3$, and $3 \to 4$ there are only $13$, $4$, and $2$ successful repairs respectively, which is too few for any contrastive direction to be estimated regardless of whether one exists. Extending the geometric analysis of Sections~\ref{subsec:raw_tests} and \ref{subsec:sensitivity_check} to later steps would require a substantially larger dataset with many more successful multi-step repairs.
